\chapter{Error-State Prediction}
\label{chap:error_state_prediction}

\section{Continuous-Time Error-State Equations}

The v1 reduced error state is
\begin{equation}
\deltax =
\begin{bmatrix}
\delta\pb^{eT} &
\delta\vb^{eT} &
\delta\boldsymbol{\theta}^{T} &
\delta\bb_g^{T} &
\delta\bb_a^{T}
\end{bmatrix}^{T}
\in\R^{15},
\end{equation}
where $\delta\boldsymbol{\theta}$ is resolved in ECEF and follows
\cref{eq:v1_attitude_error_dcm_relation}. The linearized model is
\begin{equation}
\dot{\deltax}
=
F\deltax+G\wb.
\label{eq:v1_error_state_space}
\end{equation}

Define
\begin{align}
\Omega_{ie}^{e} &= \crossmat{\omegab_{ie}^{e}},\\
\hat{\fb}_{ib}^{e} &= \hat{C}_b^e\hat{\fb}_{ib}^{b},\\
G_p^e &=
\left.\frac{\partial\gb^e}{\partial\pb_{eb}^e}\right|_{\hat{\pb}_{eb}^{e}}.
\end{align}
For Earth products, $G_p^e$ should come from the selected WGS-84/J2-capable
gravity policy rather than a hard-coded spherical approximation.

\subsection{Navigation Error Equations}

Using the convention above, the position, velocity, and attitude error
equations are grouped separately from the IMU bias equations:
\begin{align}
\delta\dot{\pb}^{e}
&=
\delta\vb^{e},\\
\delta\dot{\vb}^{e}
&=
G_p^e\delta\pb^{e}
-
2\Omega_{ie}^{e}\delta\vb^{e}
+
\crossmat{\hat{\fb}_{ib}^{e}}\delta\boldsymbol{\theta}
-
\hat{C}_b^e\delta\bb_a
+
\hat{C}_b^e\mathbf{w}_a,\\
\delta\dot{\boldsymbol{\theta}}
&=
-
\Omega_{ie}^{e}\delta\boldsymbol{\theta}
+
\hat{C}_b^e\delta\bb_g
-
\hat{C}_b^e\mathbf{w}_g.
\end{align}

\subsection{IMU Error Models}

The v1 IMU error model separates white measurement noise from bias random walk:
\begin{align}
\delta\omegab_{ib}^{b}
&\approx
-
\delta\bb_g
+
\mathbf{w}_g,\\
\delta\fb_{ib}^{b}
&\approx
-
\delta\bb_a
+
\mathbf{w}_a.
\end{align}
The corresponding IMU bias error equations are
\begin{align}
\delta\dot{\bb}_g
&=
\mathbf{w}_{bg},\\
\delta\dot{\bb}_a
&=
\mathbf{w}_{ba}.
\end{align}
These signs match the attitude-error convention in
\cref{eq:v1_attitude_error_dcm_relation} and the derivation in the broader
NavKit reference \cite{groves2013,sola2018}.

\section{Continuous-Time Matrices}

For
\begin{equation}
\wb=
\begin{bmatrix}
\mathbf{w}_g^T &
\mathbf{w}_a^T &
\mathbf{w}_{bg}^T &
\mathbf{w}_{ba}^T
\end{bmatrix}^{T},
\end{equation}
the reduced v1 continuous-time matrices are
\begin{equation}
F =
\begin{bmatrix}
\zero & \I & \zero & \zero & \zero\\
G_p^e & -2\Omega_{ie}^{e} & \crossmat{\hat{\fb}_{ib}^{e}} & \zero & -\hat{C}_b^e\\
\zero & \zero & -\Omega_{ie}^{e} & \hat{C}_b^e & \zero\\
\zero & \zero & \zero & \zero & \zero\\
\zero & \zero & \zero & \zero & \zero
\end{bmatrix},
\label{eq:v1_f_matrix}
\end{equation}
and
\begin{equation}
G =
\begin{bmatrix}
\zero & \zero & \zero & \zero\\
\zero & \hat{C}_b^e & \zero & \zero\\
-\hat{C}_b^e & \zero & \zero & \zero\\
\zero & \zero & \I & \zero\\
\zero & \zero & \zero & \I
\end{bmatrix}.
\label{eq:v1_g_matrix}
\end{equation}
Let the continuous spectral-density matrix be
\begin{equation}
\Qc =
\operatorname{diag}
\left(
S_g,\,
S_a,\,
S_{bg},\,
S_{ba}
\right).
\end{equation}
Here $S_g$ is gyro white-noise spectral density, $S_a$ is accelerometer
white-noise spectral density, and $S_{bg}$ and $S_{ba}$ are the gyro- and
accelerometer-bias random-walk spectral densities.
\begin{center}
\begin{tabular}{lll}
\toprule
Block & Meaning & Units \\
\midrule
$S_g$ & gyro white angular-rate noise PSD & $\si{\radian^2\per\second}$ \\
$S_a$ & accelerometer white specific-force noise PSD & $\si{\meter^2\per\second^3}$ \\
$S_{bg}$ & gyro-bias random-walk PSD & $\si{\radian^2\per\second^3}$ \\
$S_{ba}$ & accelerometer-bias random-walk PSD & $\si{\meter^2\per\second^5}$ \\
\bottomrule
\end{tabular}
\end{center}

The word ``reduced'' here means the v1 fifteen-state error model: position,
velocity, ECEF-resolved attitude perturbation, gyro bias, and accelerometer
bias. It does not refer to the later first-order discretization approximation.
The continuous-time equations are already a first-order linearization of the
nonlinear navigation dynamics.

\begin{warningbox}
The signs in \cref{eq:v1_f_matrix,eq:v1_g_matrix} are a convention lock. If
the implementation changes the attitude-error definition or quaternion
correction side, the prediction equations, measurement Jacobians, injection,
reset, and tests must be changed together.
\end{warningbox}

\section{Analytical Discretization}

Let $F_k$ denote the continuous-time error-state dynamics matrix from
\cref{eq:v1_f_matrix} evaluated at the nominal state for interval $k$ and held
constant over $[t_k,t_k+\Delta t]$.
With
\begin{align}
\hat{C}_{b,k}^{e} &= C_b^e\left(\hat{\quat}_{b,k}^{e}\right),\\
\hat{\fb}_{ib,k}^{e} &= \hat{C}_{b,k}^{e}\hat{\fb}_{ib,k}^{b},
\qquad
\hat{\fb}_{ib,k}^{b}
=
\frac{\Delta\mathbf{V}_{ib,k}^{b}}{\Delta t}
\quad\text{from \cref{eq:v1_interval_average_specific_force}},\\
G_{p,k}^{e} &=
\left.
\frac{\partial\gb^e}{\partial\pb_{eb}^{e}}
\right|_{\hat{\pb}_{eb,k}^{e}},
\end{align}
the v1 interval-specific continuous-time dynamics matrix is
\begin{equation}
\boxed{
F_k =
\begin{bmatrix}
\zero & \I & \zero & \zero & \zero\\
G_{p,k}^{e} & -2\Omega_{ie}^{e} & \crossmat{\hat{\fb}_{ib,k}^{e}} & \zero & -\hat{C}_{b,k}^{e}\\
\zero & \zero & -\Omega_{ie}^{e} & \hat{C}_{b,k}^{e} & \zero\\
\zero & \zero & \zero & \zero & \zero\\
\zero & \zero & \zero & \zero & \zero
\end{bmatrix}.
}
\label{eq:v1_fk_matrix}
\end{equation}

The state transition matrix (STM) over a covariance-prediction interval is the
matrix exponential of $F_k$:
\begin{equation}
\Ph_k
=
\exp(F_k\Delta t).
\end{equation}

The corresponding frozen noise-input matrix is built directly from
\cref{eq:v1_g_matrix}. The v1 interval-specific noise-input matrix is
\begin{equation}
\boxed{
G_k =
\begin{bmatrix}
\zero & \zero & \zero & \zero\\
\zero & \hat{C}_{b,k}^{e} & \zero & \zero\\
-\hat{C}_{b,k}^{e} & \zero & \zero & \zero\\
\zero & \zero & \I & \zero\\
\zero & \zero & \zero & \I
\end{bmatrix}.
}
\label{eq:v1_gk_matrix}
\end{equation}
Thus $\Ph_k$ is derived from $F_k$, while $G_k$ is evaluated directly from the
nominal attitude at interval $k$.

The general Taylor expansion for the frozen $F_k$ matrix is
\begin{equation}
\Ph_k
\approx
\I
+
F_k\Delta t
+
\frac{1}{2}F_k^2\Delta t^2
+
\frac{1}{6}F_k^3\Delta t^3
+
\cdots.
\end{equation}
The v1 baseline implementation should use the first-order state-transition
approximation
\begin{equation}
\boxed{
\Ph_k
\approx
\I+F_k\Delta t.
}
\label{eq:v1_phi_first_order}
\end{equation}

The discrete process-noise covariance must correspond to the same frozen
dynamics:
\begin{equation}
\mathbf{Q}_{d,k}
=
\int_0^{\Delta t}
\Ph_k(\tau)G_k\Qc G_k^T\Ph_k^T(\tau)d\tau.
\label{eq:v1_qd_integral}
\end{equation}
For implementation, expand the integral analytically as
\begin{equation}
\mathbf{Q}_{d,k}
\approx
\sum_{i=0}^{n}
\sum_{j=0}^{n}
\frac{
F_k^i
G_k\Qc G_k^T
\left(F_k^T\right)^j
\Delta t^{i+j+1}
}{
i!\,j!\,(i+j+1)
}.
\label{eq:v1_qd_series}
\end{equation}
This continuous-to-discrete process-noise integral and series form is the
analytical frozen-system result used in standard Kalman-filter covariance
prediction \cite{simon2006,brown2012,palumbo2010}. The v1 baseline should start
with the first-order process-noise approximation
\begin{equation}
\boxed{
\mathbf{Q}_{d,k}
\approx
G_k\Qc G_k^T\Delta t.
}
\label{eq:v1_qd_first_order}
\end{equation}
This first-order $\mathbf{Q}_{d,k}$ is the v1 implementation baseline, not a
claim that higher-order terms are physically negligible. Validation must compare
the baseline against \cref{eq:v1_qd_series} or equivalent analytical block
specializations over representative medium-rate prediction intervals. If the
first-order approximation is not adequate for the selected prediction cadence,
higher-order analytical truncation must be implemented before treating the path
as production quality.

\section{Useful First-Order Noise Blocks}

For tests and code review, the following frozen-attitude, frozen-noise blocks
are useful sanity checks. Define
$S_a^e=\hat{C}_b^eS_a\hat{C}_e^b$ and
$S_{ba}^e=\hat{C}_b^eS_{ba}\hat{C}_e^b$. If only accelerometer white noise
and accelerometer-bias random walk are considered in the
$[\delta\pb^e,\delta\vb^e,\delta\bb_a]$ substate, then
\begin{equation}
Q_{pva}
\approx
\begin{bmatrix}
\frac{\Delta t^3}{3}S_a^e+\frac{\Delta t^5}{20}S_{ba}^e &
\frac{\Delta t^2}{2}S_a^e+\frac{\Delta t^4}{8}S_{ba}^e &
-\frac{\Delta t^3}{6}\hat{C}_b^eS_{ba}\\
\frac{\Delta t^2}{2}S_a^e+\frac{\Delta t^4}{8}S_{ba}^e &
\Delta t\,S_a^e+\frac{\Delta t^3}{3}S_{ba}^e &
-\frac{\Delta t^2}{2}\hat{C}_b^eS_{ba}\\
-\frac{\Delta t^3}{6}S_{ba}\hat{C}_e^b &
-\frac{\Delta t^2}{2}S_{ba}\hat{C}_e^b &
\Delta t\,S_{ba}
\end{bmatrix}.
\end{equation}
Similarly, for the $[\delta\boldsymbol{\theta},\delta\bb_g]$ substate,
\begin{equation}
Q_{\theta g}
\approx
\begin{bmatrix}
\Delta t\,S_g^e+\frac{\Delta t^3}{3}S_{bg}^e &
\frac{\Delta t^2}{2}\hat{C}_b^eS_{bg}\\
\frac{\Delta t^2}{2}S_{bg}\hat{C}_e^b &
\Delta t\,S_{bg}
\end{bmatrix},
\end{equation}
where $S_g^e=\hat{C}_b^eS_g\hat{C}_e^b$ and
$S_{bg}^e=\hat{C}_b^eS_{bg}\hat{C}_e^b$. The full implementation should use
\cref{eq:v1_qd_series} or an equivalent analytical specialization so Earth
rate, gravity-gradient, and specific-force attitude coupling are handled
consistently.

\section{Covariance Prediction}

The covariance prediction is
\begin{equation}
\boxed{
P_{k+1}^{-}
=
\Ph_k P_k^{+}\Ph_k^T + \mathbf{Q}_{d,k}.
}
\label{eq:v1_covariance_prediction}
\end{equation}

If the Navigator buffers history for delayed or asynchronous aiding, it should
store the transition information needed to relate a measurement time to the
current filter time. The first version may keep this history short and
explicit.
